\section{本章小结}
\label{sec:push-summary}

本章开头问：原空间上的质量经可测映射后怎样分布在目标空间上，目标空间上的积分怎样回到原空间计算。第一个问题由原像回答。推前测度$T_\#\mu(B)=\mu(T^{-1}(B))$把目标集合的质量规定为它的原像的质量；原像保持不交并，所以推前测度可数可加，总质量不变，并且分步推前与一次推前相同。$\sigma$有限性则可能在推前时丢失，用到相关定理时要单独核对。

例子显示了推前测度记录什么、丢掉什么。一一映射只搬动质量；多对一映射在同一输出点汇合质量，丢掉原像之间的区分，甚至把区间上的长度汇合成点质量；伸缩改变长度的比例；坐标投影只保留柱形集合上的值，两个边缘合起来也不能恢复联合测度。

第二个问题由推前积分恒等式$\int_Yg\,\mathrm d(T_\#\mu)=\int_Xg\circ T\,\mathrm d\mu$回答。它沿积分的构造从指示函数推广到非负函数，再在绝对可积条件下推广到有符号与复值函数。恒等式只用到原像，不要求映射可逆，也不要求任何密度；它与上一章的换测度公式不同，后者点不移动，只改变累计权重。

在目标空间选定参考测度后，两步改写把原空间的积分化为目标空间的加权积分，前提是推前测度有密度；常值映射把正长度的区间压成一点，说明这个前提可以失败。

三种计算机制各有依据，不能互相替代。微分同胚的推前密度是$p(T^{-1}(\bm y))\lvert\det DT^{-1}(\bm y)\rvert$，因子在新坐标上取值、方向是逆映射，它与通常换元公式中在原坐标上取值的正向Jacobian互为对照，伸缩映射可以核对方向。多分支映射要求分支不交、例外集零测，再对落到同一点的有限个原像分支求和。坐标投影的原像是整条纤维，需要乘积参考测度，用Tonelli定理沿舍去的坐标积分；一般的非线性降维还要另行确定纤维上的参考测度，例如先用极坐标把纤维拉直再投影。

这些结论完成了测度部分的主线：从集合的测量到函数的积分，从同一空间上的换测度到空间之间的推前。第三部分将在这些积分工具之上研究函数空间的范数、完备性与逼近；流形上的体积测度与积分，即本章坐标换元在弯曲空间上的推广，在第~\ref{chap:manifold-calculus}章研究。
